% ============================================================================
% PART V: EXPERIMENTAL EVIDENCE IN FULL
% Fragment file: deepdive_part5.tex (no preamble; included from master)
% All numbers in this part are taken verbatim (or rounded) from the archived
% result files e0_structure_diagnostic.json, e1_synth.json,
% e1s_stitched_crypto.json, e2_capacity_sweep.json, e3_dormancy_sweep.json,
% e4_heterogeneity_sweep.json, e5_ablations.json, e6_snr_audit.json.
% ============================================================================

\part{Experimental Evidence in Full}
\label{sec:dd5-part}

The theory of the preceding parts makes a small number of sharp, falsifiable
claims: that a reward-independent niche assignment retains dormant-regime
competence where every reward-coupled allocator relearns; that an invariance
objective inside a niche buys robustness that an invariance objective inside a
monolith cannot; that the two mechanisms interact super-additively; and that
all of this is conditional on the environment actually containing persistent,
re-activating conditional structure. This part subjects each claim to its
experiment. We present the numbers exhaustively---every arm, every metric,
every pre-registered comparison with its raw and corrected $p$-value---and we
present the null results with the same care as the positive ones, because the
protocol's credibility rests on its ability to return ``no'' (Sections
\ref{sec:dd5-e0e1s} and \ref{sec:dd5-e6}).

\section{Methodology and Statistical Apparatus}
\label{sec:dd5-method}

\subsection{The arms}
\label{sec:dd5-arms}

Eleven arms share an identical decision layer (top-$K$ portfolio construction
over $n=20$ assets from predicted returns) and identical data per seed; they
differ only in how predictive capacity is organized and trained.

\begin{table}[h]
\centering
\caption{The eleven arms. ``Assignment'' is how specialists (or capacity
slots) are bound to regimes; ``Objective'' is the within-specialist training
loss. Oracle arms receive the true regime$\to$specialist map and serve as
skylines, not competitors.}
\label{tab:dd5-arms}
\small
\begin{tabular}{llll}
\toprule
Arm & Name & Assignment & Objective \\
\midrule
A1 & monolith-ERM & none (single model) & ERM \\
A2 & MoE router & learned gate, reward-driven & ERM \\
A3 & recent-perf allocator & capital $\propto$ recent performance & ERM \\
A4 & random fixed niches & random, frozen at init & ERM \\
A5 & RISP-ERM & learned, dormancy-pinned & ERM \\
A6 & RISP-INV (full system) & learned, dormancy-pinned & invariance \\
A7 & monolith-INV & none (single model) & invariance \\
A8a & Hedge over fixed experts & exponential weights, frozen experts & none \\
A8b & Hedge over learners & exponential weights, learning experts & ERM \\
A9 & oracle-pinned-ERM & oracle map, pinned & ERM \\
A10 & oracle-pinned-INV & oracle map, pinned & invariance \\
\bottomrule
\end{tabular}
\end{table}

The synthetic market has $n=20$ assets, $d=20$ features, $R=4$ latent regimes
with dormancy/reactivation dynamics, per-specialist capacity $K=2$, and
$20$ seeds per configuration. Evaluation is restricted to the second half of
each stream so that all arms are past their transient.

\subsection{The seeded paired design}
\label{sec:dd5-paired}

Every seed $s \in \{1,\dots,20\}$ fixes the entire data-generating draw: the
regime schedule, the per-regime parameters, the feature noise, and the
realized returns. All eleven arms are then run on the \emph{same} stream for
seed $s$. Consequently the between-seed component of variance---which regime
schedule was drawn, how harsh the reactivations happened to be---is common to
all arms, and arm differences within a seed are estimates of the treatment
effect alone. This is the experimental-design analogue of common random
numbers: it does not change the means but sharply reduces the variance of
\emph{contrasts}.

We nonetheless report unpaired Welch tests on the per-seed means
(Section~\ref{sec:dd5-welch}) rather than paired $t$-tests. This is
deliberately conservative: the paired test would have strictly more power
(the seed-level draws are positively correlated across arms), so any
comparison that survives the unpaired test would also survive the paired one,
while the converse need not hold. Where a paired analysis materially changes
the reading (the $\lambda$ ablation in Section~\ref{sec:dd5-e5}), we say so
explicitly.

\subsection{The three metrics}
\label{sec:dd5-metrics}

All metrics are means of the per-day \emph{decision regret}
\[
r_t \;=\; V^{\star}_t - V_t,
\]
where $V_t$ is the realized value of the arm's top-$K$ portfolio on day $t$
and $V^{\star}_t$ the value of the best top-$K$ portfolio in hindsight for
that day. Regret is measured in the decision space, not the prediction space,
following the smart-predict-then-optimize critique
\citep{elmachtoub2022smart}: prediction error and decision quality can
dissociate, and our research question lives entirely on the decision side.

Let $\mathcal{T}$ be the evaluation window (second half of the stream). A
\emph{reactivation event} is a day on which a regime that has been dormant
for at least $90$ days re-enters the schedule. The \emph{probe window} of an
event is its first $15$ evaluation days. Then:
\begin{align}
\texttt{post\_react} &= \text{mean of } r_t \text{ over the union of probe
windows in } \mathcal{T}, \label{eq:dd5-postreact}\\
\texttt{steady} &= \text{mean of } r_t \text{ over }
\mathcal{T} \setminus \{\text{probe windows}\}, \label{eq:dd5-steady}\\
\texttt{overall} &= \text{mean of } r_t \text{ over all of } \mathcal{T}.
\label{eq:dd5-overall}
\end{align}
\texttt{post\_react} is the primary endpoint: it isolates the days on which
retention can matter at all. \texttt{steady} measures the price paid for the
mechanism when nothing is reactivating, and \texttt{overall} aggregates both.
Each seed contributes one number per metric per arm; all inference is on
these $n=20$ per-seed means, never on the (serially dependent) daily values.

\subsection{Welch's $t$-test}
\label{sec:dd5-welch}

For arms $a$ and $b$ with per-seed means $\bar{x}_a, \bar{x}_b$, sample
variances $s_a^2, s_b^2$, and $n_a = n_b = 20$, the Welch statistic is
\begin{equation}
t \;=\; \frac{\bar{x}_a - \bar{x}_b}
{\sqrt{\dfrac{s_a^2}{n_a} + \dfrac{s_b^2}{n_b}}},
\label{eq:dd5-welch-t}
\end{equation}
referred to a $t$ distribution with the Welch--Satterthwaite degrees of
freedom
\begin{equation}
\nu \;=\;
\frac{\left(\dfrac{s_a^2}{n_a} + \dfrac{s_b^2}{n_b}\right)^{2}}
{\dfrac{(s_a^2/n_a)^2}{n_a - 1} + \dfrac{(s_b^2/n_b)^2}{n_b - 1}}.
\label{eq:dd5-welch-df}
\end{equation}
We use the unequal-variance form rather than the pooled Student $t$ because
variance heterogeneity across arms is not an edge case here---it is a
\emph{finding}. A4's defining signature is inflated between-seed variance
(its post-reactivation 95\% CI half-width is $0.00126$ against A5's
$0.00070$, a factor of $1.8$; Section~\ref{sec:dd5-e1-a4}), and A6's
high-SNR collapse comes with CI half-widths four to six times A1's
(Section~\ref{sec:dd5-e6}). Pooling variances across such pairs would
miscalibrate exactly the comparisons we care most about. With $n_a=n_b=20$
and roughly comparable variances, $\nu$ ranges over roughly $30$--$38$ in our
comparisons; the test is essentially a $z$-test with a mild small-sample
correction.

\subsection{Holm--Bonferroni and a proof that it controls FWER}
\label{sec:dd5-holm}

The E1 analysis involves a pre-registered family of $m = 11$ pairwise
comparisons (Table~\ref{tab:dd5-e1-tests}). With $11$ tests at level $0.05$
each, the chance of at least one false rejection under the global null is
bounded only by $0.55$ by Bonferroni's inequality---useless. We therefore
control the family-wise error rate (FWER) with Holm's step-down procedure,
which dominates plain Bonferroni at no assumption cost.

\begin{definition}[Holm step-down procedure]
\label{def:dd5-holm}
Order the $p$-values $p_{(1)} \le p_{(2)} \le \cdots \le p_{(m)}$ with
corresponding hypotheses $H_{(1)},\dots,H_{(m)}$. Let
\[
k^\ast = \min\Bigl\{ k : p_{(k)} > \frac{\alpha}{m - k + 1} \Bigr\}
\]
(with $k^\ast = m+1$ if no such $k$ exists). Reject
$H_{(1)},\dots,H_{(k^\ast - 1)}$ and retain the rest. Equivalently, report
adjusted $p$-values
$\tilde{p}_{(k)} = \max_{j \le k}\, \min\{(m-j+1)\, p_{(j)},\, 1\}$
and reject where $\tilde{p} \le \alpha$. The \texttt{holm\_p} entries in our
result files are these adjusted values.
\end{definition}

\begin{proposition}[Holm controls FWER at level $\alpha$]
\label{prop:dd5-holm-fwer}
Under arbitrary dependence among the $p$-values, the probability that Holm's
procedure rejects at least one true null hypothesis is at most $\alpha$.
\end{proposition}

\begin{proof}
Let $I_0$ be the index set of true nulls and $m_0 = |I_0|$; assume $m_0 \ge 1$
(otherwise there is nothing to prove). Suppose the procedure rejects at least
one true null, and let $k$ be the \emph{first} step (in the sorted order) at
which a true null is rejected. All hypotheses rejected at steps
$1,\dots,k-1$ are then false nulls, of which there are $m - m_0$ in total;
hence $k - 1 \le m - m_0$, i.e.\ $k \le m - m_0 + 1$, and therefore
\[
\frac{\alpha}{m - k + 1} \;\le\; \frac{\alpha}{m_0}.
\]
Because step $k$ rejected, $p_{(k)} \le \alpha/(m-k+1) \le \alpha/m_0$, and
since the true null rejected at step $k$ has $p$-value $p_{(k)}$,
\[
\min_{i \in I_0} p_i \;\le\; p_{(k)} \;\le\; \frac{\alpha}{m_0}.
\]
So the event ``at least one true null is rejected'' implies the event
$\{\min_{i \in I_0} p_i \le \alpha/m_0\}$. For each true null,
$\Pr(p_i \le \alpha/m_0) \le \alpha/m_0$ (validity of the individual test),
so by the union bound
\[
\Pr\Bigl(\min_{i \in I_0} p_i \le \frac{\alpha}{m_0}\Bigr)
\;\le\; m_0 \cdot \frac{\alpha}{m_0} \;=\; \alpha. \qedhere
\]
\end{proof}

Multiple-testing discipline is not optional in this literature: the data
mining of trading-strategy families is precisely the failure mode documented
by the reality-check and SPA literature
\citep{white2000reality,hansen2005spa}, by the deflated Sharpe ratio
\citep{bailey2014deflated}, and by the backtesting critiques of
\citet{harvey2016backtesting} and \citet{lopez2018advances}. Our family is
small ($11$), fixed in advance, and corrected; that is the minimum standard
those papers demand.

\subsection{Confidence intervals and the choice of 20 seeds}
\label{sec:dd5-ci}

All reported intervals are per-arm 95\% confidence intervals on the mean over
seeds,
\begin{equation}
\bar{x} \;\pm\; t_{19,\,0.975}\, \frac{s}{\sqrt{20}}, \qquad
t_{19,\,0.975} = 2.093,
\label{eq:dd5-ci}
\end{equation}
where $s$ is the between-seed standard deviation. The \texttt{ci95} fields in
the result files are the half-widths of these intervals.

Why $20$ seeds? Back out the per-seed SD from Eq.~\eqref{eq:dd5-ci}: on the
primary endpoint the major arms have $s \approx 0.0014$--$0.0027$ (A1:
$0.0020$; A5: $0.0015$; A6: $0.0014$; A4: $0.0027$). The tightest
pre-registered contrast is the invariance axis A6 vs.\ A5, with observed
difference $0.00142$. The Welch standard error for that pair is
$\sqrt{(0.0015^2 + 0.0014^2)/20} \approx 0.00046$, giving $t \approx 3.1$ and
the observed raw $p = 0.0061$---comfortably detectable, with power
$\approx 0.85$ at $\alpha = 0.05$ for an effect of that size. Halving the
seed count would have left the invariance axis at the edge of detectability
after Holm correction; doubling it would have bought little, since the
binding constraint on the smallest effects we care about (e.g.\ A6 vs.\ A10)
is that they are \emph{predicted to be null}. Twenty seeds is the point where
every effect we predicted to exist is powered and every effect we predicted
not to exist is bounded by a usefully narrow CI.

\subsection{The noise-floor convention}
\label{sec:dd5-floor}

Raw regret levels conflate two things: irreducible per-day noise (even the
oracle-assigned, invariance-trained A10 cannot drive regret to zero in a
stochastic market) and the mechanism-dependent excess above that floor. We
adopt the convention
\begin{equation}
\text{floor} \;=\; \texttt{steady}(\text{A10}) \;=\; 0.0250, \qquad
\text{excess}(\cdot) \;=\; \text{metric}(\cdot) - 0.0250.
\label{eq:dd5-floor}
\end{equation}
A10's steady state is the best any arm in this family achieves with oracle
assignment, oracle-free training, and no reactivation stress; treating it as
the irreducible floor makes relative improvements interpretable. A $10.3\%$
improvement in raw regret (A6 vs.\ A1, post-reactivation) is a $38.1\%$
reduction in excess-over-floor---the latter is the honest measure of how much
of the \emph{addressable} regret the mechanism removes
(Section~\ref{sec:dd5-e1-axes}).

\subsection{Pre-registration discipline}
\label{sec:dd5-prereg}

Before the final 20-seed E1 run, we fixed in writing: (i) the eleven-arm
roster and all hyperparameters; (ii) the three metrics and the primary
endpoint (\texttt{post\_react}); (iii) the family of eleven pairwise
comparisons and the Holm correction over exactly that family; (iv) a
directional prediction for every comparison, including four predicted
\emph{nulls} (A2$\approx$A1, A7$\approx$A1, A4$\approx$A5,
A6$\approx$A10); and (v) the E0 structure gate for real data, with the
commitment to run E1s and report it regardless of the gate's outcome. The
predicted nulls matter as much as the predicted effects: a theory that only
predicts differences can survive almost any data, whereas predicting that
invariance alone does \emph{nothing} post-reactivation (A7$\approx$A1) is a
claim the data could easily have destroyed.

\begin{intuition}
The entire E1 design is a $2\times 2$ factorial hiding inside an eleven-arm
lineup: \{monolith, niche\} $\times$ \{ERM, INV\} is the square
(A1, A5, A7, A6), and everything else is a control that pins down \emph{why}
the square comes out the way it does. A2 and A3 show that reward-coupled
allocation does not reproduce the niche effect; A4 shows assignment quality
matters beyond mere partitioning; A8a/A8b bracket the system between
``memory without adaptation'' and ``adaptation without memory''; A9/A10 show
how close learned pinning gets to the oracle ceiling.
\end{intuition}

\begin{warning}
\textbf{What these statistics do not establish.} Every $p$-value in this
part is a statement about \emph{mean decision regret on a specified
generator or historical sample, under a specified decision layer, relative
to other arms of the same family}. None of the following is claimed, and the
reader should resist inferring any of it: (i) no claim of positive risk-adjusted
return (``alpha'') on any market---regret relative to a daily hindsight
oracle is not profit; (ii) no claim about live trading, transaction costs,
slippage, capacity, or market impact; (iii) no claim that the synthetic
generator is a calibrated model of any asset class; (iv) no claim that the
real-data nulls of Section~\ref{sec:dd5-e0e1s} would persist under other
universes, frequencies, or feature sets---they are scoped to the data
actually tested. The expected post-publication fate of any genuine
market anomaly is decay \citep{mclean2016does}; our claims are deliberately
about \emph{mechanism}, which is the only thing a synthetic-plus-null-real
evidence base can support.
\end{warning}

% ============================================================================
\section{E1: The Headline Dissociation, Number by Number}
\label{sec:dd5-e1}

\subsection{The full table}
\label{sec:dd5-e1-table}

Table~\ref{tab:dd5-e1-main} reports all eleven arms on all three metrics,
with 95\% CIs, from the archived 20-seed run ($K=2$, heterogeneity $1.0$,
hard memory, probe $15$, minimum dormancy $90$ days).

\begin{table}[h]
\centering
\caption{E1 headline results. Mean daily decision regret ($\times 10^{-2}$,
lower is better) $\pm$ 95\% CI half-width over 20 seeds. Excess = post-react
mean minus the noise floor $2.50 \times 10^{-2}$
(Eq.~\eqref{eq:dd5-floor}). Bold: best non-oracle arm per column.}
\label{tab:dd5-e1-main}
\small
\begin{tabular}{lccccc}
\toprule
Arm & \texttt{post\_react} & \texttt{overall} & \texttt{steady} &
Excess & Post/steady \\
\midrule
A1 monolith-ERM      & $3.424 \pm 0.095$ & $2.690 \pm 0.044$ & $2.459 \pm 0.058$ & $0.925$ & $1.39$ \\
A2 router            & $3.408 \pm 0.110$ & $2.695 \pm 0.052$ & $2.481 \pm 0.070$ & $0.909$ & $1.37$ \\
A3 recent-perf       & $3.776 \pm 0.096$ & $2.940 \pm 0.055$ & $2.685 \pm 0.072$ & $1.277$ & $1.41$ \\
A4 random-fixed      & $3.295 \pm 0.126$ & $2.827 \pm 0.114$ & $2.717 \pm 0.124$ & $0.795$ & $1.21$ \\
A5 RISP-ERM      & $3.214 \pm 0.070$ & $2.750 \pm 0.064$ & $2.644 \pm 0.084$ & $0.715$ & $1.22$ \\
A6 RISP-INV      & $\mathbf{3.072 \pm 0.066}$ & $\mathbf{2.602 \pm 0.059}$ & $2.507 \pm 0.074$ & $\mathbf{0.572}$ & $1.23$ \\
A7 monolith-INV      & $3.435 \pm 0.100$ & $2.663 \pm 0.044$ & $\mathbf{2.428 \pm 0.062}$ & $0.935$ & $1.41$ \\
A8a Hedge-fixed      & $4.390 \pm 0.090$ & $3.815 \pm 0.069$ & $3.704 \pm 0.096$ & $1.890$ & $1.19$ \\
A8b Hedge-learners   & $3.407 \pm 0.100$ & $2.672 \pm 0.046$ & $2.443 \pm 0.057$ & $0.907$ & $1.39$ \\
\midrule
A9 oracle-pinned-ERM & $3.210 \pm 0.077$ & $2.736 \pm 0.066$ & $2.624 \pm 0.088$ & $0.710$ & $1.22$ \\
A10 oracle-pinned-INV& $3.053 \pm 0.078$ & $2.590 \pm 0.059$ & $2.499 \pm 0.075$ & $0.553$ & $1.22$ \\
\bottomrule
\end{tabular}
\end{table}

Three structural facts are visible before any test is run. First, the
\texttt{post\_react} column separates the arms far more than the
\texttt{steady} column: retention is a property of reactivation days.
Second, the post/steady ratio splits the arms into two clean families:
everything with a reward-coupled or absent assignment (A1, A2, A3, A7, A8b)
sits at $1.37$--$1.41$, and everything with a fixed or pinned assignment
(A4, A5, A6, A8a, A9, A10) sits at $1.19$--$1.23$. The ratio is a
fingerprint of the mechanism, independent of the level. Third, the
\texttt{steady} column shows the honest cost: pinning is not free. A5's
steady regret ($2.644$) is \emph{worse} than A1's ($2.459$), because
partitioned specialists see less data per niche; A7 actually posts the best
steady value in the table ($2.428$), better even than A10. RISP wins
\texttt{overall} because reactivation days are frequent and expensive enough
that retention pays for the partitioning cost---not because it dominates
everywhere.

\subsection{The pre-registered comparisons}
\label{sec:dd5-e1-tests}

\begin{table}[h]
\centering
\caption{The eleven pre-registered comparisons on \texttt{post\_react}.
$\Delta$ = first arm minus second arm ($\times 10^{-2}$; negative favors the
first arm). Verdicts at FWER $\alpha = 0.05$ (Holm). ``Null (pred.)''
denotes comparisons pre-registered as approximate equalities.}
\label{tab:dd5-e1-tests}
\small
\begin{tabular}{llrccl}
\toprule
\# & Comparison & $\Delta$ & Welch $p$ & Holm $p$ & Verdict \\
\midrule
1 & A6 vs A1  & $-0.353$ & $9.2\times 10^{-7}$ & $1.0\times 10^{-5}$ & rejected: A6 $<$ A1 \\
2 & A6 vs A7  & $-0.363$ & $1.2\times 10^{-6}$ & $1.2\times 10^{-5}$ & rejected: A6 $<$ A7 \\
3 & A6 vs A8b & $-0.335$ & $4.5\times 10^{-6}$ & $4.1\times 10^{-5}$ & rejected: A6 $<$ A8b \\
4 & A5 vs A1  & $-0.210$ & $1.3\times 10^{-3}$ & $1.0\times 10^{-2}$ & rejected: A5 $<$ A1 \\
5 & A6 vs A5  & $-0.142$ & $6.1\times 10^{-3}$ & $4.3\times 10^{-2}$ & rejected: A6 $<$ A5 \\
6 & A5 vs A2  & $-0.194$ & $6.3\times 10^{-3}$ & $4.3\times 10^{-2}$ & rejected: A5 $<$ A2 \\
7 & A6 vs A9  & $-0.138$ & $1.1\times 10^{-2}$ & $5.7\times 10^{-2}$ & \emph{marginal} (fails Holm) \\
8 & A5 vs A4  & $-0.081$ & $0.281$ & $1.0$ & null (pred.) \\
9 & A6 vs A10 & $+0.019$ & $0.719$ & $1.0$ & null (pred.) \\
10 & A2 vs A1 & $-0.016$ & $0.827$ & $1.0$ & null (pred.) \\
11 & A7 vs A1 & $+0.011$ & $0.882$ & $1.0$ & null (pred.) \\
\bottomrule
\end{tabular}
\end{table}

We now walk through each pre-registered prediction.

\paragraph{P1: A6 $<$ A5 $<$ A1 (both axes contribute).}
Confirmed at every link. A5 beats A1 by $0.210 \times 10^{-2}$
(Holm $p = 0.010$): pinned niches alone, with plain ERM inside them, retain
dormant-regime competence. A6 beats A5 by a further $0.142 \times 10^{-2}$
(Holm $p = 0.043$): the invariance objective adds value \emph{on top of}
retention. The chain A6 $<$ A1 is the largest effect in the family
(Holm $p = 1.0 \times 10^{-5}$).

\paragraph{P2: A7 $\approx$ A1 post-reactivation (the dissociation).}
Confirmed as a null: $\Delta = +0.011 \times 10^{-2}$, raw $p = 0.882$. The
invariance objective, trained in a monolith, does \emph{nothing} for
post-reactivation regret---indeed A7's point estimate is fractionally worse
than A1's. This is the single most theory-laden prediction in the family:
invariance can only protect parameters that still exist, and a monolith
overwrites its dormant-regime parameters regardless of how they were
trained. Note the contrast with \texttt{steady}, where A7 is the best arm in
the table: invariance helps the monolith \emph{between} reactivations and is
helpless \emph{at} them.

\paragraph{P3: A2 $\approx$ A1 (a reward-driven router is a monolith in
disguise).}
Confirmed as a null: $\Delta = -0.016 \times 10^{-2}$, raw $p = 0.827$. The
learned gate receives no gradient from a dormant regime and reallocates its
experts' capacity to live regimes; on reactivation it relearns like the
monolith. The companion test A5 $<$ A2 (Holm $p = 0.043$) confirms the
positive side: dormancy-pinning beats the router by about as much as it
beats the monolith.

\paragraph{P4: A3 is the worst learner.}
Confirmed. A3's post-reactivation regret ($3.776 \times 10^{-2}$) is the
worst of any learning arm---only the frozen A8a is worse---and uniquely
among the learners its damage is \emph{not} confined to probe windows: its
\texttt{overall} ($2.940$) and \texttt{steady} ($2.685$) are also the worst
in the learner family. The failure mode is continuous erosion: the
recent-performance allocator drives capital toward whatever just worked,
and the capital floor (the minimum allocation that keeps a specialist
training at all) is exactly where a temporarily-unlucky specialist's data
diet collapses. Every regime transition triggers a starve-and-relearn cycle,
so A3 pays the relearning tax perpetually, not just after long dormancies.

\paragraph{P5: A8a retains but cannot adapt; A8b adapts but cannot retain.}
Confirmed on both jaws of the bracket. A8a (Hedge over frozen experts) has
the \emph{lowest} post/steady ratio in the table ($1.19$): its relative
degradation at reactivation is minimal, because nothing it could forget was
ever learned online---this is retention in its purest form. But its absolute
level is catastrophic ($\texttt{steady} = 3.704$, $51\%$ above A1), because
frozen experts cannot track within-regime drift at all. A8b (Hedge over
learning experts) is the mirror image: its level tracks the monolith almost
exactly ($\texttt{post\_react}$ $3.407$ vs $3.424$; \texttt{steady} $2.443$
vs $2.459$; ratio $1.39$ vs $1.39$), because exponential weights over
experts that all share the reward stream do not protect any expert's
parameters from being retrained. The sleeping-experts story, in other words,
is not about the \emph{weighting} scheme; it is about whether the experts'
parameters are insulated from reward-coupled reuse. A6 beats A8b by
$0.335\times 10^{-2}$ (Holm $p = 4.1\times 10^{-5}$).

\paragraph{P6: A4 $\approx$ A5 in mean, with higher variance.}
\label{sec:dd5-e1-a4}
Confirmed. Random fixed niches land at $3.295 \times 10^{-2}$ against A5's
$3.214$ (raw $p = 0.281$, Holm $1.0$): a random reward-independent
assignment captures most of the retention effect, exactly as the theory
says it should---retention comes from reward-independence, not from
assignment quality. But A4's between-seed SD is $0.0027$ against A5's
$0.0015$ ($1.8\times$), and the signature is even starker on
\texttt{steady} (CI $0.124$ vs $0.084$) and \texttt{overall} (CI $0.114$ vs
$0.064$). Inspection of A4's per-seed values shows why: most seeds are fine,
but the seeds in which the random assignment leaves a regime with no
well-matched specialist (coverage gaps) blow up the mean (per-seed
\texttt{overall} values up to $0.0341$ against a median near $0.0270$).
Learned pinning buys insurance against the coverage tail, not a better
average draw.

\paragraph{P7: A6 $\approx$ A10 (the skyline is reached).}
Confirmed as a null: $\Delta = +0.019 \times 10^{-2}$, raw $p = 0.719$. The
full system is statistically indistinguishable from the same system with an
oracle-given assignment. Learned dormancy-pinning recovers essentially all
of the oracle assignment's value; the corresponding ERM-side check is even
cleaner (A5 $3.214$ vs A9 $3.210$, not in the test family but visibly
identical).

\paragraph{P8 (exploratory edge): A6 vs A9.}
A6's point estimate beats the \emph{oracle-assigned ERM} arm by
$0.138\times 10^{-2}$---learned assignment plus invariance edging out
perfect assignment without it. Raw $p = 0.011$, but the Holm-adjusted value
is $0.057$, which fails the $0.05$ FWER threshold. We report this exactly
as it is: suggestive, consistent with the rest of the pattern, and not
established at the pre-registered standard. Nothing downstream depends on
it.

\begin{keypoint}
Seven of seven predicted differences went the predicted direction; six of
seven survived Holm at $\alpha=0.05$ (the seventh, A6 $<$ A9, is marginal
at $0.057$). All four predicted nulls came back null with raw
$p \ge 0.28$. No comparison in the pre-registered family contradicted a
prediction.
\end{keypoint}

\subsection{The axes arithmetic}
\label{sec:dd5-e1-axes}

Decomposing the headline effect along the factorial square:
\begin{align}
\text{retention axis (A1} \to \text{A5):}\quad
& \frac{3.424 - 3.214}{3.424} = 6.1\%, \label{eq:dd5-axis-ret}\\[2pt]
\text{invariance axis (A5} \to \text{A6):}\quad
& \frac{3.214 - 3.072}{3.214} = 4.4\%, \label{eq:dd5-axis-inv}\\[2pt]
\text{joint (A1} \to \text{A6):}\quad
& \frac{3.424 - 3.072}{3.424} = 10.3\%. \label{eq:dd5-axis-joint}
\end{align}
In excess-over-floor terms (Eq.~\eqref{eq:dd5-floor}) the joint effect is
\begin{equation}
1 - \frac{3.072 - 2.500}{3.424 - 2.500}
= 1 - \frac{0.572}{0.925} = 38.1\%:
\label{eq:dd5-excess}
\end{equation}
the full system removes a bit over a third of the addressable
post-reactivation regret, and lands within $3\%$ of the oracle skyline's
excess ($0.553$ vs $0.572$).

\subsection{The interaction index: super-additivity}
\label{sec:dd5-e1-interaction}

Define, per seed, the factorial interaction
\begin{equation}
I \;=\; (\text{A1} - \text{A5}) + (\text{A1} - \text{A7})
- (\text{A1} - \text{A6})
\;=\; \text{A1} - \text{A5} - \text{A7} + \text{A6}.
\label{eq:dd5-interaction}
\end{equation}
If the two mechanisms contributed additively, $I$ would be zero: the joint
gain would equal the sum of the marginal gains. Computed per seed on the
paired design (here pairing is legitimate and free, since
Eq.~\eqref{eq:dd5-interaction} is a within-seed contrast),
\[
I \;=\; -0.0015 \;\pm\; 0.0005 \quad (\text{mean} \pm \text{95\% CI},\;
n = 20),
\]
which excludes zero decisively. The sign is the interesting part: $I < 0$
means the joint gain $(\text{A1}-\text{A6}) = 0.0035$ \emph{exceeds} the sum
of the marginals $(\text{A1}-\text{A5}) + (\text{A1}-\text{A7}) =
0.0021 + (-0.0001) = 0.0020$. The interaction is super-additive, and the
arithmetic shows why in one line: the invariance marginal in the monolith is
\emph{zero} (A1 $-$ A7 $= -0.0001$), yet inside niches invariance
contributes $0.0014$ (A5 $-$ A6). Invariance is not an independent additive
bonus; it is a multiplier on retained memories. A specialist that still
holds its dormant regime's parameters can hold \emph{invariant} parameters,
which transfer across the within-regime drift accumulated during dormancy;
a monolith has nothing for the invariance penalty to protect at
reactivation time. This is the empirical signature of the
retention$\times$robustness interaction derived in
Part~III, and it echoes the
invariance-for-decision-making rationale of \citet{krueger2021rex} and
\citet{yan2026invpnco}: invariance pays where the deployment distribution
differs from the training distribution---which, for a retained dormant
niche, is exactly the reactivation moment.

% ============================================================================
\section{E0 and E1s: The Real-Data Nulls}
\label{sec:dd5-e0e1s}

A mechanism paper earns the right to a synthetic headline only by trying,
honestly, to find the mechanism's preconditions in real data, and reporting
what it finds. We found that the precondition---persistent, re-activating,
\emph{decision-relevant} conditional structure---is absent in the real
universes we tested. This section documents that null in full.

\subsection{E0: the structure diagnostic}
\label{sec:dd5-e0}

\paragraph{Protocol.} E0 asks the gating question without running any
learning arm at all: \emph{if an oracle could condition the decision policy
on regime labels, would that beat not conditioning?} The components:

\begin{enumerate}
\item \textbf{Universes.} Five crypto pairs (daily bars) and three
commodities (daily). Evaluation spans $562$ days (crypto) and $1{,}270$
days (commodities).
\item \textbf{Causal labelers.} Two regime labelers, L1 (finer) and L2
(coarser), computed from \emph{trailing} data only, so labels are known
before the day's decision and cannot leak the target. The realized regime
distributions are imbalanced, especially under L2 on crypto
(counts $517/517/31/59$), which we note as a power limitation for the rare
regimes.
\item \textbf{Decision metric.} The same top-$K$ portfolio regret as E1
(Section~\ref{sec:dd5-metrics})---a decision-space test, not a
prediction-space test.
\item \textbf{Oracle conditioning.} Fit the best per-regime decision policy
with hindsight model selection per regime (an upper bound on what any
regime-aware learner could extract), and compare its mean regret to
(a) the same machinery run on \emph{block-shuffled} labels (preserving label
marginals and block-length statistics, destroying the label--day
alignment), and (b) a pooled, unconditioned policy.
\item \textbf{Statistic.} Gap $\% = (\text{shuffled} - \text{real}) /
\text{shuffled} \times 100$ (positive iff real labels help), and
$z = (\text{shuffled} - \text{real}) / \mathrm{sd}(\text{shuffled})$ over
the shuffle distribution.
\end{enumerate}

\begin{table}[h]
\centering
\caption{E0 structure diagnostic on real data. Mean daily decision regret
($\times 10^{-2}$) for oracle regime-conditioning on real labels, on
block-shuffled labels (mean $\pm$ SD of the shuffle distribution), and for
the pooled unconditioned policy. Positive gap = real labels help. No cell
shows exploitable structure.}
\label{tab:dd5-e0}
\small
\begin{tabular}{llccrrr}
\toprule
Domain & Labeler & Real & Shuffled & Pooled & Gap (\%) & $z$ \\
\midrule
Crypto ($n{=}562$ d) & L1 & $1.220$ & $1.202 \pm 0.027$ & $1.195$ & $-1.47$ & $-0.66$ \\
Crypto ($n{=}562$ d) & L2 & $1.212$ & $1.204 \pm 0.047$ & $1.195$ & $-0.66$ & $-0.17$ \\
Commodities ($n{=}1270$ d) & L1 & $0.834$ & $0.832 \pm 0.015$ & $0.816$ & $-0.25$ & $-0.14$ \\
Commodities ($n{=}1270$ d) & L2 & $0.836$ & $0.838 \pm 0.009$ & $0.816$ & $+0.16$ & $+0.16$ \\
\bottomrule
\end{tabular}
\end{table}

\paragraph{Reading.} Across both domains and both labelers, the real-label
oracle never beats the shuffled-label oracle by more than $0.16\%$, all
$|z| < 0.7$, and three of four point estimates are \emph{negative}---real
labels do marginally worse than shuffled ones. Most damning, the
\emph{pooled} policy is the best cell in every row: even hindsight-optimal
regime conditioning loses to not conditioning at all, because splitting the
sample by regime costs more in estimation variance than the (apparently
nonexistent) regime structure returns. This is a null with teeth: it bounds
not just RISP but \emph{any} regime-conditioned method in this family,
on this data, at this frequency, with these features.

\paragraph{Honest scoping.} The null is scoped: five crypto pairs and three
commodities, daily bars, one feature set, two specific (if reasonable)
causal labelers, and a top-$K$ decision layer. Intraday frequencies, larger
cross-sections, alternative labelers (e.g.\ macro-indexed), or other
decision layers could in principle revive the gate, and the protocol is
designed to be rerun cheaply on any of them.

\paragraph{Relation to the GAUSE null.} The predecessor project found a
closely related null in \emph{prediction} space: regime-aware prediction
showed no real-vs-shuffled gap on financial domains \citep{li2026gause}.
That did not settle our question, and pretending otherwise would have been
sloppy: the predict-then-optimize literature documents real cases where
prediction-space and decision-space rankings dissociate
\citep{elmachtoub2022smart,yan2026invpnco}, so a decision-metric retest was
a genuine open question, not a formality. E0 answers it: the dissociation
does not rescue financial regime structure here. The prediction-space null
and the decision-space null now agree.

\subsection{E1s: the stitched-crypto experiment}
\label{sec:dd5-e1s}

E0 was the pre-registered gate for real-data claims, and it failed.
Pre-registration committed us to running the full arm lineup on real data
anyway and predicting \emph{null separation}---a falsifiable prediction in
its own right (if arms had separated despite the failed gate, the gate
would be broken). E1s stitches real crypto regime blocks (block counts
$15/13/7/9$ across the four labeled regimes) into schedules with genuine
dormancy (minimum $80$ days, probe $15$ days), $20$ seeds of stitching
randomization, ten arms (A10 requires generator access and is undefined on
real data).

\begin{table}[h]
\centering
\caption{E1s stitched real crypto. Mean daily decision regret
($\times 10^{-2}$) $\pm$ 95\% CI, 20 seeds. All pairwise comparisons in the
pre-registered family: raw Welch $p > 0.21$, all Holm $p = 1.0$.}
\label{tab:dd5-e1s}
\small
\begin{tabular}{lcc}
\toprule
Arm & \texttt{post\_react} & \texttt{overall} \\
\midrule
A1 monolith-ERM      & $1.396 \pm 0.051$ & $1.285 \pm 0.017$ \\
A2 router            & $1.403 \pm 0.050$ & $1.296 \pm 0.020$ \\
A3 recent-perf       & $1.398 \pm 0.075$ & $1.299 \pm 0.025$ \\
A4 random-fixed      & $1.373 \pm 0.057$ & $1.273 \pm 0.027$ \\
A5 RISP-ERM      & $1.381 \pm 0.056$ & $1.276 \pm 0.025$ \\
A6 RISP-INV      & $1.344 \pm 0.064$ & $1.263 \pm 0.025$ \\
A7 monolith-INV      & $1.410 \pm 0.057$ & $1.290 \pm 0.018$ \\
A8a Hedge-fixed      & $1.325 \pm 0.053$ & $1.239 \pm 0.023$ \\
A8b Hedge-learners   & $1.400 \pm 0.059$ & $1.285 \pm 0.017$ \\
A9 oracle-pinned-ERM & $1.342 \pm 0.052$ & $1.268 \pm 0.027$ \\
\bottomrule
\end{tabular}
\end{table}

All ten arms land in a band of $1.32$--$1.41 \times 10^{-2}$
post-reactivation; the six pre-registered comparisons return raw
$p \in \{0.22, 0.23, 0.40, 0.57, 0.71, 0.95\}$ and Holm $p = 1.0$
throughout. Two details confirm that this is structure-absence rather than
mere noise. First, the numerically best post-reactivation arm is
\emph{A8a}---the frozen-expert arm that cannot adapt to anything---which is
exactly what one expects when there is nothing regime-specific to retain or
relearn: not adapting is as good as adapting. Second, A9 (oracle labels!)
is statistically indistinguishable from A6 with learned labels
($p = 0.95$): even perfect knowledge of the regime labels confers nothing,
because the labels do not predict decision-relevant structure (which is
what E0 measured directly).

\begin{keypoint}
\textbf{The protocol is falsifiable, and falsifies.} The identical
machinery---same arms, same metrics, same test family, same Holm
correction---produces Holm $p \approx 10^{-5}$ separations where the
generator contains dormant-but-persistent structure (E1), and a uniform
$p > 0.2$ flatline where a pre-registered diagnostic says the structure is
absent (E0 $\to$ E1s). A pipeline that can only return positives is
untrustworthy on synthetic data too; the E1s flatline is what makes the E1
separation credible.
\end{keypoint}

% ============================================================================
\section{E2--E4: Boundary Sweeps}
\label{sec:dd5-sweeps}

The headline lives at one point in configuration space ($K=2$, dormancy
$\ge 90$, heterogeneity $1.0$). E2--E4 sweep the three axes the theory says
should govern the effect: capacity pressure, dormancy length, and
cross-regime heterogeneity.

\subsection{E2: capacity}
\label{sec:dd5-e2}

The retention argument requires capacity pressure: a monolith forgets a
dormant regime only if serving the live regimes forces reuse of its
capacity. As $K \to R$ the pressure vanishes and the retention gap should
close. E2 sweeps $K \in \{1,2,3,4\}$ ($R = 4$) under two memory models:
\emph{hard} (slot eviction) and \emph{soft} (graded interference decay).

\begin{table}[h]
\centering
\caption{E2 capacity sweep, \texttt{post\_react} ($\times 10^{-2}$) $\pm$
95\% CI, hard memory. The pinned arms (A5, A6, A9) are exactly constant
across $K$: each specialist serves only its own niches, so the capacity
constraint never binds for them. Rightmost column: retention gap
A1 $-$ A5 ($\times 10^{-2}$).}
\label{tab:dd5-e2-hard}
\small
\begin{tabular}{cccccccc}
\toprule
$K$ & A1 & A2 & A4 & A5 & A6 & A9 & A1$-$A5 \\
\midrule
1 & $3.427 \pm 0.102$ & $3.432 \pm 0.100$ & $3.667 \pm 0.163$ & $3.214 \pm 0.070$ & $3.072 \pm 0.066$ & $3.210 \pm 0.077$ & $0.213$ \\
2 & $3.424 \pm 0.095$ & $3.408 \pm 0.110$ & $3.295 \pm 0.126$ & $3.214 \pm 0.070$ & $3.072 \pm 0.066$ & $3.210 \pm 0.077$ & $0.210$ \\
3 & $3.303 \pm 0.088$ & $3.276 \pm 0.064$ & $3.202 \pm 0.072$ & $3.214 \pm 0.070$ & $3.072 \pm 0.066$ & $3.210 \pm 0.077$ & $0.089$ \\
4 & $3.221 \pm 0.082$ & $3.228 \pm 0.068$ & $3.206 \pm 0.078$ & $3.214 \pm 0.070$ & $3.072 \pm 0.066$ & $3.210 \pm 0.077$ & $0.007$ \\
\bottomrule
\end{tabular}
\end{table}

\begin{table}[h]
\centering
\caption{E2 capacity sweep, \texttt{post\_react} ($\times 10^{-2}$) $\pm$
95\% CI, soft memory (no eviction; graded interference with
overflow-scaled decay). Rightmost column: retention gap A1 $-$ A5.}
\label{tab:dd5-e2-soft}
\small
\begin{tabular}{cccccccc}
\toprule
$K$ & A1 & A2 & A4 & A5 & A6 & A9 & A1$-$A5 \\
\midrule
1 & $3.276 \pm 0.079$ & $3.285 \pm 0.067$ & $3.667 \pm 0.163$ & $3.214 \pm 0.070$ & $3.072 \pm 0.066$ & $3.210 \pm 0.077$ & $0.062$ \\
2 & $3.277 \pm 0.079$ & $3.254 \pm 0.077$ & $3.318 \pm 0.117$ & $3.214 \pm 0.070$ & $3.072 \pm 0.066$ & $3.210 \pm 0.077$ & $0.063$ \\
3 & $3.279 \pm 0.079$ & $3.265 \pm 0.056$ & $3.239 \pm 0.071$ & $3.214 \pm 0.070$ & $3.072 \pm 0.066$ & $3.210 \pm 0.077$ & $0.065$ \\
4 & $3.277 \pm 0.081$ & $3.255 \pm 0.064$ & $3.273 \pm 0.079$ & $3.214 \pm 0.070$ & $3.072 \pm 0.066$ & $3.210 \pm 0.077$ & $0.063$ \\
\bottomrule
\end{tabular}
\end{table}

Four observations.

\begin{enumerate}
\item \textbf{Reward-independent arms are exactly flat.} A5, A6, and A9 post
identical values at every $K$ in both memory models. This is not a rounding
artifact but a structural fact: a pinned specialist serves only its assigned
niches, so its memory never overflows and the capacity model is inert for
it. Flatness is the cleanest possible demonstration that the pinned arms'
advantage is \emph{not} extra effective capacity---they use the same budget
and simply never churn it.
\item \textbf{Gap closure at $K = R$ under hard memory.} The retention gap
A1 $-$ A5 falls from $0.213$ at $K=1$ to $0.007$ at $K=4$---zero to within
a tenth of a CI. With enough slots, the monolith never evicts, never
forgets, and pinning buys nothing. The invariance edge survives, however:
A6 ($3.072$) still beats everything at $K=4$, confirming the axes are
distinct (capacity governs retention, not robustness).
\item \textbf{Non-closure under soft memory---at any capacity.} With no
eviction, capacity slots stop binding and the reward-driven arms go
\emph{flat} in $K$ at $\approx 3.28$: gentler than hard eviction at tight
capacity ($3.43$ at $K \le 2$) but never closing the gap to the pinned
arms ($\approx 0.063$ above A5, $0.21$ above A6, at every $K$), because
graded interference from foreign updates erodes dormant heads regardless
of how many slots exist. The hard model predicts closure at $K = R$ and
closes; the soft model predicts a capacity-independent residue and shows
exactly that.
\item \textbf{A4's coverage problem is worst at $K=1$.} Random fixed niches
post $3.667 \pm 0.163$ at $K=1$---the worst learner cell in the
sweep---because with one slot per specialist a single unlucky assignment
leaves a regime entirely uncovered. The cell's CI ($0.163$) is the widest in
the sweep, the variance signature of Section~\ref{sec:dd5-e1-a4} amplified.
\end{enumerate}

\subsection{E3: dormancy length}
\label{sec:dd5-e3}

Retention should matter more the longer a regime sleeps: the monolith's
relearning debt accumulates with dormancy, while a pinned specialist's
retained parameters do not decay (hard memory). E3 sweeps the minimum
dormancy threshold $D \in \{21, 63, 126, 252, 504\}$ days.

\begin{table}[h]
\centering
\caption{E3 dormancy sweep ($\times 10^{-2}$, $\pm$ 95\% CI). Each $D$ is a
\emph{different schedule generator} (shorter dormancies mean more frequent
transitions and different steady-state mixes), so levels are comparable
\emph{within} a row block only, never across $D$. Rightmost columns:
within-$D$ retention$+$invariance gap A1 $-$ A6, absolute and relative.}
\label{tab:dd5-e3}
\small
\begin{tabular}{clccr@{\quad}r}
\toprule
$D$ & Arm & \texttt{post\_react} & \texttt{steady} &
\multicolumn{2}{c}{A1 $-$ A6 (post)} \\
\midrule
\multirow{5}{*}{21}
& A1 & $5.226 \pm 0.148$ & $3.225 \pm 0.225$ & & \\
& A2 & $5.241 \pm 0.151$ & $3.305 \pm 0.213$ & & \\
& A3 & $5.475 \pm 0.169$ & $3.504 \pm 0.263$ & $0.152$ & $2.9\%$ \\
& A5 & $5.198 \pm 0.167$ & $3.416 \pm 0.220$ & & \\
& A6 & $5.074 \pm 0.142$ & $3.293 \pm 0.222$ & & \\
\midrule
\multirow{5}{*}{63}
& A1 & $5.024 \pm 0.202$ & $2.719 \pm 0.214$ & & \\
& A2 & $5.028 \pm 0.208$ & $2.785 \pm 0.226$ & & \\
& A3 & $5.362 \pm 0.173$ & $2.973 \pm 0.218$ & $0.375$ & $7.5\%$ \\
& A5 & $4.818 \pm 0.198$ & $2.876 \pm 0.208$ & & \\
& A6 & $4.648 \pm 0.193$ & $2.776 \pm 0.209$ & & \\
\midrule
\multirow{5}{*}{126}
& A1 & $4.975 \pm 0.160$ & $2.550 \pm 0.157$ & & \\
& A2 & $4.960 \pm 0.150$ & $2.598 \pm 0.156$ & & \\
& A3 & $5.228 \pm 0.147$ & $2.747 \pm 0.159$ & $0.401$ & $8.1\%$ \\
& A5 & $4.787 \pm 0.185$ & $2.762 \pm 0.141$ & & \\
& A6 & $4.574 \pm 0.181$ & $2.649 \pm 0.130$ & & \\
\midrule
\multirow{5}{*}{252}
& A1 & $4.336 \pm 0.098$ & $2.633 \pm 0.100$ & & \\
& A2 & $4.344 \pm 0.103$ & $2.657 \pm 0.096$ & & \\
& A3 & $4.682 \pm 0.103$ & $2.811 \pm 0.103$ & $0.385$ & $8.9\%$ \\
& A5 & $4.095 \pm 0.144$ & $2.777 \pm 0.099$ & & \\
& A6 & $3.951 \pm 0.103$ & $2.652 \pm 0.092$ & & \\
\midrule
\multirow{5}{*}{504}
& A1 & $3.668 \pm 0.059$ & $2.530 \pm 0.080$ & & \\
& A2 & $3.646 \pm 0.059$ & $2.565 \pm 0.080$ & & \\
& A3 & $3.968 \pm 0.049$ & $2.713 \pm 0.079$ & $0.408$ & $11.1\%$ \\
& A5 & $3.454 \pm 0.076$ & $2.668 \pm 0.077$ & & \\
& A6 & $3.260 \pm 0.061$ & $2.533 \pm 0.077$ & & \\
\bottomrule
\end{tabular}
\end{table}

\begin{warning}
Do not compare levels across $D$ blocks. A $D = 21$ generator produces a
chaotic schedule in which \emph{every} arm is perpetually post-transition
(note the elevated steady values, $3.2$--$3.5$, against $2.5$--$2.7$ at
$D = 504$); a $D = 504$ generator produces long calm stretches. The only
licensed reading is the \emph{within-$D$ gap between arms}, which holds the
schedule fixed.
\end{warning}

The within-$D$ gap A1 $-$ A6 grows steeply from $0.152$ at $D = 21$ to
$0.375$ at $D = 63$, then saturates in absolute terms
($0.375 \to 0.401 \to 0.385 \to 0.408$), while the relative gap continues
to drift up ($7.5\% \to 11.1\%$) because the baseline level falls as
schedules calm down. The shape is exactly the theory's: at $D = 21$ the
monolith's parameters have barely drifted from the dormant regime, so there
is little to retain; by $D \approx 63$--$126$ (one to two quarters of
dormancy) the monolith's relearning debt is fully formed and further
dormancy adds little, because the debt is bounded by ``relearn from
scratch.'' The ordering A6 $<$ A5 $<$ A1 $\approx$ A2, and A3 worst, holds
at every $D$ without exception.

\subsection{E4: heterogeneity}
\label{sec:dd5-e4}

The invariance axis should engage only insofar as regimes differ: at
heterogeneity $0$ all regimes share one mechanism and there is nothing to be
invariant \emph{across}. E4 sweeps the cross-regime heterogeneity scale
$h \in \{0, 0.5, 1.0, 1.5, 2.0\}$ for the pinned arms.

\begin{table}[h]
\centering
\caption{E4 heterogeneity sweep, \texttt{post\_react} ($\times 10^{-2}$)
$\pm$ 95\% CI. Rightmost column: ERM$-$INV gap (A5 $-$ A6).}
\label{tab:dd5-e4}
\small
\begin{tabular}{ccccc}
\toprule
$h$ & A5 RISP-ERM & A6 RISP-INV & A9 oracle-pinned-ERM & A5$-$A6 \\
\midrule
0.0 & $2.998 \pm 0.080$ & $2.900 \pm 0.079$ & $2.997 \pm 0.081$ & $0.097$ \\
0.5 & $3.073 \pm 0.081$ & $2.948 \pm 0.080$ & $3.067 \pm 0.073$ & $0.125$ \\
1.0 & $3.214 \pm 0.070$ & $3.072 \pm 0.066$ & $3.210 \pm 0.077$ & $0.142$ \\
1.5 & $3.459 \pm 0.087$ & $3.249 \pm 0.095$ & $3.418 \pm 0.094$ & $0.210$ \\
2.0 & $3.749 \pm 0.138$ & $3.511 \pm 0.124$ & $3.679 \pm 0.165$ & $0.238$ \\
\bottomrule
\end{tabular}
\end{table}

Three readings.

\begin{enumerate}
\item \textbf{Both arms degrade with heterogeneity, ERM faster.} Fitting
endpoint slopes, A5 degrades at $0.376 \times 10^{-2}$ per unit $h$ against
A6's $0.306$---the ERM arm is roughly $1.2\times$ more sensitive to
heterogeneity than the INV arm, and the oracle-pinned ERM arm A9 sits in
between ($0.341$). The cleaner statement is in the gap itself: the
ERM$-$INV gap grows \emph{monotonically} and roughly $2.4$-fold across the
sweep, from $0.0010$ at $h = 0$ to $0.0024$ at $h = 2$
(in raw units, $0.00097 \to 0.00238$). The invariance objective buys more
the more the regimes differ---which is what an invariance objective is for.
\item \textbf{The $h = 0$ residual gap.} At zero heterogeneity the theory's
first-order prediction is gap $= 0$, yet A6 retains a small, CI-overlapping
but consistently signed edge ($0.097 \times 10^{-2}$). The likely mechanism
is mundane and worth naming: with $d = 20$ features of which only a subset
is causal, the cross-environment variance penalty also suppresses reliance
on \emph{spuriously correlated distractor features within} a niche---it
acts as a targeted regularizer even when there is no cross-regime
heterogeneity to be invariant to. We flag this as a confound on the pure
invariance interpretation at low $h$: part of the A5$-$A6 gap at every $h$
is plain regularization, and only the \emph{growth} of the gap with $h$
(the $0.0010 \to 0.0024$ monotone trend) is unambiguously attributable to
invariance proper.
\item \textbf{A5 tracks A9 everywhere.} Learned pinning matches
oracle pinning under ERM at every heterogeneity ($\le 0.041 \times 10^{-2}$
apart, well inside CIs), extending the P7 skyline result along the whole
sweep.
\end{enumerate}

% ============================================================================
\section{E5--E6: Ablations and the Inversion Audit}
\label{sec:dd5-e5e6}

\subsection{E5: ablations of the full system}
\label{sec:dd5-e5}

E5 ablates A6's four design choices one at a time, holding everything else
at the headline configuration; the response variable is
\texttt{post\_react} ($\times 10^{-2}$, 20 seeds).

\paragraph{The invariance weight $\beta$.}

\begin{table}[h]
\centering
\caption{E5 $\beta$ sweep (A6 variant, \texttt{post\_react},
$\times 10^{-2}$). Reference points: A5 retained-ERM $= 3.214 \pm 0.070$;
headline A6 ($\beta = 1$) $= 3.072 \pm 0.066$.}
\label{tab:dd5-e5-beta}
\small
\begin{tabular}{cccccc}
\toprule
$\beta$ & $0$ & $0.25$ & $1$ & $4$ & $16$ \\
\midrule
\texttt{post\_react} & $3.555 \pm 0.107$ & $3.090 \pm 0.079$ &
$3.072 \pm 0.066$ & $3.170 \pm 0.076$ & $3.522 \pm 0.065$ \\
\bottomrule
\end{tabular}
\end{table}

The response is a clean U-shape with minimum at $\beta = 1$ and a broad
plateau over $\beta \in [0.25, 4]$ (all within $\sim 3\%$ of the minimum,
CIs overlapping). Both extremes collapse, and the collapses are
informative:
\begin{itemize}
\item At $\beta = 0$ the within-niche objective degenerates to the
cross-environment variance term alone, and performance ($3.555$) falls not
merely to the retained-ERM level but \emph{below} it
(A5 $= 3.214$; the variance-only arm is $10.6\%$ worse). This is precisely
the degeneracy the theory predicts (and that the V-REx analysis warns of
\citep{krueger2021rex}): a variance-only criterion is minimized by any
predictor that \emph{equalizes} risk across environments, including
uniformly bad constant-like predictors. Equal risks are not low risks. The
empirical collapse below the ERM reference is the theory's degenerate
solution showing up on schedule.
\item At $\beta = 16$ the risk term swamps the variance term and
performance ($3.522$) gives back the entire invariance gain and more,
landing $14.6\%$ above the $\beta = 1$ optimum and visibly worse than plain
retained ERM. We read the overshoot beyond A5 cautiously (the heavily
re-weighted objective also rescales effective step sizes), but the
qualitative point stands: the benefit requires \emph{both} terms in
balance, and a factor-of-16 window around the default ($0.25$--$4$) is all
flat---the choice $\beta = 1$ is not a tuned knife-edge.
\end{itemize}

\paragraph{The environment window $W_c$.}

\begin{table}[h]
\centering
\caption{E5 $W_c$ sweep (\texttt{post\_react}, $\times 10^{-2}$).}
\label{tab:dd5-e5-wc}
\small
\begin{tabular}{ccccc}
\toprule
$W_c$ & $1$ & $5$ & $20$ & $60$ \\
\midrule
\texttt{post\_react} & $3.033 \pm 0.076$ & $3.057 \pm 0.069$ &
$3.072 \pm 0.066$ & $3.090 \pm 0.070$ \\
\bottomrule
\end{tabular}
\end{table}

We had entertained a deflationary explanation of the invariance gain: that
chunking each niche's stream into $W_c$-day pseudo-environments effectively
\emph{batches} the updates, and the gain is a batching/smoothing artifact
rather than an invariance effect. The sweep refutes this cleanly, and we
report it as the refutation it is: across a $60\times$ range of $W_c$ the
response is flat to within $0.057 \times 10^{-2}$---less than one CI
half-width---and the numerically best cell is $W_c = 1$, i.e.\ \emph{no}
batching at all. If batching were the active ingredient, $W_c = 1$ would be
the worst cell, not the best. The pseudo-environment construction matters
for defining the variance penalty, not for any smoothing side effect.

\paragraph{The assignment-smoothing parameter $\lambda$.}
With $\lambda$ off ($0$): $3.141 \pm 0.082$; on ($0.3$): $3.072 \pm 0.066$.
The unpaired CIs overlap, so by the conservative standard of
Section~\ref{sec:dd5-paired} the effect is not individually established.
The paired within-seed contrast, which is legitimate here, gives
$+0.069 \pm 0.060$ ($\times 10^{-2}$, mean $\pm$ 95\% CI)---a $2.2\%$
improvement whose interval excludes zero, barely. Honest summary:
$\lambda$ helps a little, with weak evidence; nothing in the headline
depends on it.

\paragraph{Pinning schedule: pinned vs.\ never.}
Removing dormancy-pinning entirely (assignments stay perpetually plastic)
degrades \texttt{post\_react} from $3.072 \pm 0.066$ to $3.197 \pm 0.092$;
the paired contrast is $+0.126 \pm 0.059$ ($\times 10^{-2}$), a $3.9\%$
loss that is clearly resolved. Never-pinning gives back most of the
invariance-axis gain (the never-pin arm sits near the A5/A9 retained-ERM
level of $3.21$), which is consistent with the interaction story of
Section~\ref{sec:dd5-e1-interaction}: a perpetually plastic assignment lets
reactivation-time reassignment churn destroy exactly the retained-invariant
parameters the mechanism exists to protect.

\subsection{E6: the SNR audit and the inversion}
\label{sec:dd5-e6}

\paragraph{Development history, stated plainly.} This experiment exists
because our first prototype failed. The initial generator was, we later
understood, calibrated to an effectively high signal-to-noise ratio, and
the arm lineup showed no positive separation---RISP did not beat the
monolith. The tempting path was to tune the generator until the predicted
effect appeared and report that point. Instead we made SNR an explicit
axis, swept it, and report the whole curve, including the region where our
method \emph{loses}. The audit below is the honest version of what could
otherwise have been a silent hyperparameter search.

\begin{table}[h]
\centering
\caption{E6 SNR audit ($\times 10^{-2}$, $\pm$ 95\% CI, 20 seeds per cell).
$\rho$ = post/steady relearning-penalty ratio. ``A6 adv.'' =
$(\text{A1}-\text{A6})/\text{A1}$ on \texttt{post\_react}; positive favors
RISP.}
\label{tab:dd5-e6}
\scriptsize
\begin{tabular}{clccc r}
\toprule
SNR & Arm & \texttt{post\_react} & \texttt{steady} & $\rho$ & A6/A9 adv. \\
\midrule
\multirow{3}{*}{$0.5\times$}
& A1  & $3.778 \pm 0.077$ & $2.963 \pm 0.057$ & $1.28$ & \\
& A6  & $3.466 \pm 0.073$ & $2.890 \pm 0.066$ & $1.20$ & $+8.3\%$ \\
& A9  & $3.584 \pm 0.068$ & $3.000 \pm 0.068$ & $1.19$ & $+5.1\%$ \\
\midrule
\multirow{3}{*}{$1\times$}
& A1  & $3.424 \pm 0.095$ & $2.459 \pm 0.058$ & $1.39$ & \\
& A6  & $3.072 \pm 0.066$ & $2.507 \pm 0.074$ & $1.23$ & $+10.3\%$ \\
& A9  & $3.210 \pm 0.077$ & $2.624 \pm 0.088$ & $1.22$ & $+6.3\%$ \\
\midrule
\multirow{3}{*}{$2\times$}
& A1  & $2.990 \pm 0.075$ & $1.878 \pm 0.062$ & $1.59$ & \\
& A6  & $2.779 \pm 0.117$ & $2.273 \pm 0.110$ & $1.22$ & $+7.0\%$ \\
& A9  & $2.908 \pm 0.121$ & $2.369 \pm 0.128$ & $1.23$ & $+2.7\%$ \\
\midrule
\multirow{3}{*}{$4\times$}
& A1  & $2.733 \pm 0.070$ & $1.480 \pm 0.091$ & $1.85$ & \\
& A6  & $3.089 \pm 0.265$ & $2.577 \pm 0.217$ & $1.20$ & $-13.1\%$ \\
& A9  & $3.210 \pm 0.283$ & $2.660 \pm 0.250$ & $1.21$ & $-17.5\%$ \\
\midrule
\multirow{3}{*}{$8\times$}
& A1  & $3.286 \pm 0.155$ & $1.594 \pm 0.210$ & $2.06$ & \\
& A6  & $4.905 \pm 0.620$ & $4.117 \pm 0.465$ & $1.19$ & $-49.3\%$ \\
& A9  & $4.939 \pm 0.629$ & $4.194 \pm 0.535$ & $1.18$ & $-50.3\%$ \\
\bottomrule
\end{tabular}
\end{table}

\paragraph{The inversion.} The RISP advantage on the primary endpoint
is $+8.3\%$ at $0.5\times$, peaks at $+10.3\%$ at the headline $1\times$,
narrows to $+7.0\%$ at $2\times$, then \emph{inverts}: $-13.1\%$ at
$4\times$ and $-49.3\%$ at $8\times$. The crossover sits between $2\times$
and $4\times$, i.e.\ around $2$--$3\times$ the headline SNR.

\paragraph{The relearning-penalty ratio isolates the mechanism.} The
post/steady ratio $\rho$ tells the mechanistic story with unusual clarity.
A6's ratio is \emph{flat} across the entire sweep:
$1.20, 1.23, 1.22, 1.20, 1.19$. The retention mechanism works identically
at every SNR---a pinned specialist's reactivation penalty is always the
same modest $\sim\!20\%$ premium over its own steady state. A1's ratio, by
contrast, climbs monotonically: $1.28, 1.39, 1.59, 1.85, 2.06$. The
monolith's \emph{relative} relearning penalty gets worse with SNR, not
better. So why does the monolith win at high SNR? Because its
\emph{absolute levels} collapse: A1's steady regret falls from $2.96$ at
$0.5\times$ to $1.48$ at $4\times$, while A6's steady regret
\emph{deteriorates} ($2.89 \to 2.58 \to 4.12$ from $0.5\times$ through
$8\times$, non-monotonically). At high SNR a pooled learner extracts the
strong signal from all data jointly and relearns a reactivated regime
within days---faster than the probe window can penalize it---whereas each
pinned specialist is confined to its own niche's fraction of the data and
carries a variance penalty calibrated for a noisier world. Pooling beats
partitioning when the signal is strong enough to relearn cheaply. (A1's
own non-monotonicity at $8\times$---post-react $3.29$ vs $2.73$ at
$4\times$, with steady CIs widening---suggests additional dynamics at
extreme SNR, e.g.\ aggressive-fit instability, that we have not dissected;
we note it rather than interpret it.)

\paragraph{The skyline inverts identically.} The decisive control: A9, the
\emph{oracle-assigned} pinned arm, inverts in lockstep with A6 ($-17.5\%$
at $4\times$, $-50.3\%$ at $8\times$; their CIs overlap heavily at both
points). The inversion therefore cannot be an artifact of RISP's
learned assignment, its pinning heuristic, or any estimation failure in
our mechanism---perfect assignment inverts too. The inversion is a property
of the \emph{regime}: partitioned, reward-independent memory is the right
architecture in low-SNR, slow-relearning environments and the wrong one in
high-SNR, fast-relearning environments. The boundary is part of the claim,
not a caveat appended to it.

\begin{keypoint}
\textbf{Deployment-boundary reading.} The audit gives an operational rule:
RISP-style partitioned retention should be deployed only where the
post-reactivation relearning cost of a pooled learner is demonstrably
large---empirically, below roughly $2$--$3\times$ the headline SNR in this
family, i.e.\ in environments where a monolith's post/steady ratio is
materially elevated ($\rho \gtrsim 1.4$). The monolith's measured $\rho$
is itself an observable diagnostic: it can be estimated on any stream
\emph{before} committing to the architecture. Financial daily data, where
the E0 diagnostic already found no exploitable conditional structure at
all, sits firmly outside the deployment region; the mechanism's natural
habitat is low-SNR streams with long-dormancy recurring structure.
\end{keypoint}

\begin{intuition}
A useful way to compress E6: \emph{retention is a hedge against expensive
relearning}. The price of the hedge is the partitioning cost (less data per
specialist, a variance penalty, a worse steady state). The value of the
hedge is the relearning debt it cancels. Low SNR makes relearning slow and
expensive, so the hedge is cheap relative to its payout; high SNR makes
relearning nearly free, so the hedge is all premium and no payout. The
post/steady ratios are precisely the two sides of this ledger, read off
the data.
\end{intuition}

% ============================================================================
\section{Threats to Validity}
\label{sec:dd5-threats}

We enumerate the limits of the evidence, in roughly decreasing order of
importance, with what each would take to overturn.

\begin{enumerate}
\item \textbf{The headline is synthetic, and the real-data result is null.}
This is the central limitation and we state it without hedging: every
positive number in Sections \ref{sec:dd5-e1}--\ref{sec:dd5-e5e6} comes from
a generator we wrote. E0/E1s show that the mechanism's precondition is
absent in the real universes we tested, so the contribution is a
\emph{mechanism, a boundary map, and a falsifiable deployment gate}---not a
demonstrated real-market improvement. A reader who needs the latter should
treat this paper as a protocol to run on their own data, starting with E0.
\item \textbf{Small real universes.} The E0/E1s null rests on five crypto
pairs and three commodities at daily frequency with one feature set and two
labelers. The null is well-measured on that data ($|z| < 0.7$ across all
four cells) but could fail to generalize to intraday frequencies, equities
cross-sections, or richer conditioning information. The gate is cheap to
rerun; we make no claim beyond the cells in Table~\ref{tab:dd5-e0}.
\item \textbf{Single decision-layer family.} All arms share one decision
layer (top-$K$ portfolio from predicted returns). The decision-space
dissociation literature \citep{elmachtoub2022smart,yan2026invpnco} cuts
both ways: just as prediction nulls did not settle the decision question,
our decision-layer-specific results do not settle other decision layers
(e.g.\ mean-variance with estimated covariances, or execution-aware
objectives). The within-family rankings are clean; the family is one
point in a larger space.
\item \textbf{Linear specialists.} Specialists are linear in the $d = 20$
features. Retention-vs-relearning economics could shift with
high-capacity nonlinear specialists, in either direction (faster relearning
hurts retention's value, as E6 shows; worse interference helps it). The
GAUSE neural-expert results \citep{li2026gause} suggest the
reward-independence principle transfers to gradient-trained experts, but
we have not verified that here.
\item \textbf{Task-incremental labels are given.} Arms receive the regime
label stream (synthetic: true labels; E1s: causal labeler output). The
harder class-incremental problem---detecting reactivations without
labels---is not addressed, and the GAUSE real-data experience (label-free
recovery collapsing under feature overlap) warns against assuming it is
easy.
\item \textbf{Hyperparameter choices.} Each load-bearing choice is ablated
where it matters: $\beta$ has a $16\times$ flat plateau
(Table~\ref{tab:dd5-e5-beta}), $W_c$ is flat over $60\times$
(Table~\ref{tab:dd5-e5-wc}), $\lambda$ is nearly irrelevant, pinning
matters and is the mechanism itself. The probe length ($15$ days) and
dormancy threshold ($90$ days) were fixed before the final run and the E3
sweep shows the effect is not knife-edged in the dormancy axis; we did not
sweep probe length, and the post/steady split depends on it mechanically.
\item \textbf{Seed count and inference.} $n = 20$ seeds powers the
pre-registered effects (Section~\ref{sec:dd5-ci}) but leaves the marginal
comparison A6 vs A9 unresolved (Holm $p = 0.057$) and gives only moderate
resolution on variance signatures (A4). The conservative unpaired tests
mean reported significances are, if anything, understated.
\item \textbf{What would change our conclusions.} Concretely: (i) an E0 run
on some real universe showing a robust positive gap, followed by an E1s-style
run in which arms \emph{fail} to separate, would break the protocol's
central link; (ii) a demonstration that A2-style routers with explicit
idle-expert protection match A6 would collapse the
reward-independence claim to an implementation detail; (iii) an E6-style
audit on a low-SNR regime showing the monolith's $\rho$ elevated but no
RISP advantage would falsify the deployment gate. All three are
runnable with the released code, and the third is the one we would run
first.
\end{enumerate}

\begin{keypoint}
\textbf{Summary of Part V.} On a generator with the hypothesized structure:
retention is worth $6.1\%$, invariance $4.4\%$, jointly $10.3\%$
($38.1\%$ of excess-over-floor), with a super-additive interaction
($I = -0.0015 \pm 0.0005$) and a reached skyline (A6 $\approx$ A10,
$p = 0.72$); the effect requires capacity pressure (gap $\to 0.007$ at
$K = R$, hard memory), grows then saturates with dormancy, scales with
heterogeneity (gap $0.0010 \to 0.0024$), and survives every ablation with
flat plateaus. On real daily crypto/commodities the precondition is absent
(all E0 $|z| < 0.7$) and the arms correctly fail to separate (all
E1s $p > 0.2$). At high SNR the advantage inverts ($-49.3\%$ at $8\times$),
the oracle skyline inverts with it, and the monolith's own post/steady
ratio provides a measurable deployment gate. The mechanism is real,
bounded, and honest about where it lives.
\end{keypoint}
